\documentclass[10pt]{beamer}

\usepackage{appendixnumberbeamer}
\usepackage{booktabs}
\usepackage[scale=2]{ccicons}
\usepackage{pgfplots}
\usepackage{xspace}
\usepackage{bookmark}
\usepackage{amssymb}
\usepackage{mathtools}
\usepackage[normalem]{ulem}
\usepackage[T1]{fontenc}
\usepackage[sfdefault,book]{FiraSans}
\usepackage{FiraMono}
\usepackage{fontawesome}
\usepackage{listings}
\usepackage{hyperref}
\usepackage[10pt]{moresize}

\renewcommand*\familydefault{\ttdefault}

\renewcommand{\thefootnote}{\fnsymbol{footnote}}
\renewcommand{\MintedPygmentize}{/home/thiago/.local/bin/pygmentize}

\definecolor{light-gray}{gray}{0.95}
\renewcommand\lstlistingname{Código}
\DeclareCaptionFormat{listing} {
  \parbox{\textwidth}{\hspace{-0.2cm}#1#2#3}
}
\DeclareCaptionFont{black}{\color{black}}
\captionsetup[lstlisting]{
  format=listing,
  labelfont=black,
  textfont=black,
  singlelinecheck=true,
  margin=0pt,
  font={tt,footnotesize,bf}
}

\lstset{
  basicstyle=\footnotesize\ttfamily,
  escapeinside={\%*}{*)},
  mathescape=true,
  showspaces=false,
  showtabs=false,
  showstringspaces=false,%
  % backgroundcolor=\color{light-gray},
  rulesepcolor=\color{black},
  frame=shadowbox,
  literate=
  {á}{{\'a}}1 {é}{{\'e}}1 {í}{{\'i}}1 {ó}{{\'o}}1 {ú}{{\'u}}1
  {Á}{{\'A}}1 {É}{{\'E}}1 {Í}{{\'I}}1 {Ó}{{\'O}}1 {Ú}{{\'U}}1
  {à}{{\`a}}1 {è}{{\`e}}1 {ì}{{\`i}}1 {ò}{{\`o}}1 {ù}{{\`u}}1
  {À}{{\`A}}1 {È}{{\'E}}1 {Ì}{{\`I}}1 {Ò}{{\`O}}1 {Ù}{{\`U}}1
  {ä}{{\"a}}1 {ë}{{\"e}}1 {ï}{{\"i}}1 {ö}{{\"o}}1 {ü}{{\"u}}1
  {Ä}{{\"A}}1 {Ë}{{\"E}}1 {Ï}{{\"I}}1 {Ö}{{\"O}}1 {Ü}{{\"U}}1
  {â}{{\^a}}1 {ê}{{\^e}}1 {î}{{\^i}}1 {ô}{{\^o}}1 {û}{{\^u}}1
  {Â}{{\^A}}1 {Ê}{{\^E}}1 {Î}{{\^I}}1 {Ô}{{\^O}}1 {Û}{{\^U}}1
  {Ã}{{\~A}}1 {ã}{{\~a}}1 {Õ}{{\~O}}1 {õ}{{\~o}}1
  {œ}{{\oe}}1 {Œ}{{\OE}}1 {æ}{{\ae}}1 {Æ}{{\AE}}1 {ß}{{\ss}}1
  {ű}{{\H{u}}}1 {Ű}{{\H{U}}}1 {ő}{{\H{o}}}1 {Ő}{{\H{O}}}1
  {ç}{{\c c}}1 {Ç}{{\c C}}1 {ø}{{\o}}1 {å}{{\r a}}1 {Å}{{\r A}}1
  {€}{{\euro}}1 {£}{{\pounds}}1 {«}{{\guillemotleft}}1
  {»}{{\guillemotright}}1 {ñ}{{\~n}}1 {Ñ}{{\~N}}1 {¿}{{?`}}1
}

\setminted{
  fontsize=\footnotesize,
  style=xcode,
  frame=single,
  framesep=3\fboxsep,
  labelposition=topline,
}

\title{Análise e Projeto de Algoritmos}
\author{Thiago Cavalcante  -- thiago.cavalcante@penedo.ufal.br}
\institute{Universidade Federal de Alagoas -- UFAL \\ Campus Arapiraca \\ Unidade de Ensino de Penedo}
\titlegraphic{\hfill\includegraphics[height=1.5cm]{images/brasao-ufal.eps}}


\subtitle{Aula 1}
\date{25 de outubro de 2019}

\begin{document}

\maketitle

\begin{frame}[fragile]{Informações Gerais da Disciplina}
    \begin{center}
        \begin{tabular}{@{}ll@{}}
            \toprule
            \textbf{Nome} & Análise e Projeto de Algoritmos \\
            \textbf{Código} & SISB014 \\
            \textbf{Semestre} & 5º \\
            \textbf{Carga Horária} & 80h \\
            \textbf{PPC} & 03/2018 \\
            \textbf{Turma} & 2019.2 \\
            \textbf{Horário} & Sexta, 18:10 -- 21:40 \\
            \textbf{Sala} & 4 \\
            \textbf{Lista de discussão} & Google Groups: \href{https://groups.google.com/forum/#!forum/sisb014_20192}{\alert{sisb014\_20192}} \\
            \textbf{Repositório} & GitHub: \href{https://github.com/theagoliveira/sisb014_20192}{\alert{theagoliveira/sisb014\_20192}} \\
            \bottomrule
        \end{tabular}
    \end{center}
\end{frame}

\begin{frame}{Cronograma}
    \begin{center}
        \begin{tabular}{@{}r@{/}r@{/}rll@{}}
            \toprule
            \multicolumn{4}{@{}l}{\textbf{Datas importantes}} \\
            \midrule
            22 & 11 & 2019 & \textit{Possível} dia da \alert{AB1} \\
            23 & 11 & 2019 & Prazo final para digitação da \alert{AB1} \\
            14 & 02 & 2020 & \textit{Possível} dia da \alert{AB2} \\
            17 & 02 & 2020 & Prazo final para digitação da \alert{AB2} \\
            17--22 & 02 & 2020 & Período de \alert{reavaliação} \\
            27--29 & 02 & 2020 & Período de \alert{provas finais} \\
            \bottomrule
        \end{tabular}
    \end{center}
\end{frame}

\begin{frame}{Ementa}
    \begin{itemize}
        \item Conceitos básicos: motivação e solução de problemas; critérios de análise; correção e eficiência
        \item Análise de complexidade de algs.: custo; tempo de processamento; operações elementares; função de complexidade; classes de problemas; comparação de algs.
        \item Comportamento assintótico de funções: dominação assintótica; big O; theta; omega
        \item Técnicas de análise de algs.: equações de recorrência; teorema mestre
        \item Paradigmas de projeto de algs.: indução matemática; recursividade; algs. tentativa e erro; divisão e conquista; balanceamento; programação dinâmica; algs. gulosos e aproximados
        \item Análise de algs. de busca e ordenação; grafos; conhecidos
        \item Tratamento de problemas NP-completos
    \end{itemize}
\end{frame}

\begin{frame}{Bibliografia}
    \huge
    SKIENA, S. S. The Algorithm Design Manual
\end{frame}

\begin{frame}{Introdução ao Design de Algoritmos}
    \huge
    \metroset{block=fill}
    \begin{block}{\huge Algoritmo}
        Procedimento para realizar uma determinada tarefa
    \end{block}

    \bigskip

    \begin{block}{\huge Programa de computador}
        Algoritmo + estruturas de dados
    \end{block}
\end{frame}

\begin{frame}{Introdução ao Design de Algoritmos}
    \huge
    Um algoritmo resolve um \textbf{problema}. O problema é especificado pela \textbf{instância}.

    \bigskip

    problema $\neq$ instância
\end{frame}

\begin{frame}{Introdução ao Design de Algoritmos}
    \Large

    \textbf{Exemplo}: Ordenação

    \textbf{Entrada}: Uma sequência de $n$ elementos $a_1, \ldots, a_n$

    \textbf{Saída}: A permutação (reordenação) da sequência de entrada de forma que $a^{'}_{1} \leq a^{'}_{2} \leq \cdots \leq a^{'}_{n-1} \leq a^{'}_{n}$
\end{frame}

\begin{frame}{Introdução ao Design de Algoritmos}
    \Large

    \textbf{Instâncias}

    \textbf{Entrada 1}: \{``João'', ``Thiago'', ``Maria'', ``Ana''\}

    \textbf{Entrada 2}: \{100, 34, 57, 943, 48\}

    O algoritmo transforma qualquer possível instância na saída desejada.
\end{frame}

\begin{frame}{Introdução ao Design de Algoritmos}
    \huge
    \textbf{Exemplo}: \textit{insertion sort}
\end{frame}

\begin{frame}{Introdução ao Design de Algoritmos}
    \huge

    Três propriedades de um bom algoritmo:

    \begin{itemize}
        \only<1>{\item \huge Correto}
        \only<2>{\item \huge \alert{Correto}}
        \item \huge Eficiente
        \item \huge Fácil de implementar
    \end{itemize}
\end{frame}

\begin{frame}{Introdução ao Design de Algoritmos}
    \huge
    Correção nem sempre é óbvia!

    \onslide<2->{Correção requer uma \alert{prova}.}
\end{frame}

\begin{frame}{Introdução ao Design de Algoritmos}
    \huge
    \textbf{Exemplo}: Otimização de caminho de um robô

    \Large
    \textbf{Entrada}: Conjunto de pontos em um plano

    \textbf{Saída}: Menor caminho fechado que passa por todos os pontos
\end{frame}

\begin{frame}[fragile]{Introdução ao Design de Algoritmos}
    \Large Heurística do vizinho mais próximo
    \normalsize
    \begin{lstlisting}
VizinhoMaisPróximo($P$)
  Escolha um ponto inicial $p0$ de $P$ e o visite
  $p$ = $p_0$
  $i$ = 0
  Enquanto existirem pontos não visitados
    $i = i + 1$
    Seja $p_i$ o ponto mais próximo não visitado de $p_{i-1}$
    Visite $p_i$
  Retorne para $p_0$ a partir de $p_{n-1}$
    \end{lstlisting}
\end{frame}

\begin{frame}[fragile]{Introdução ao Design de Algoritmos}
    \Large Heurística do par mais próximo
    \normalsize
    \begin{lstlisting}
ParMaisPróximo($P$)
  Seja $n$ o número de pontos no conjunto $P$
  Para $i = 1$ até $n - 1$ faça
    $d = \infty$
    Para cada par $(s, t)$ de pontas em cadeias distintas
        Se $dist(s, t) \leq d$, então $s_m = s$ e $t_m = t$ e $d = dist(s, t)$
    Conecte $(s_m, t_m)$
  Conecte as duas pontas restantes
    \end{lstlisting}
\end{frame}

\begin{frame}[fragile]{Introdução ao Design de Algoritmos}
    \begin{lstlisting}
PCVÓtimo($P$)
  $d = \infty$
  Para cada $n!$ permutações $P_i$ do conjunto $P$
    Se $(custo(P_i) \leq d)$, então $d = custo(P_i)$ e $P_{min} = P_i$
  Retorne $P_{min}$
    \end{lstlisting}

    \onslide<2->{\large 20 pontos: só 2.432.902.008.176.640.000 permutações}

    \bigskip

    \onslide<3->{\large 1000 pontos: ??????}

    \bigskip

    \onslide<3->{\large \textbf{PCV}: Problema do Caixeiro-viajante}
\end{frame}

\begin{frame}{Introdução ao Design de Algoritmos}
    \huge
    \textbf{Algoritmos} \alert{sempre} fornecem o resultado correto, \textbf{heurísticas} fazem um bom trabalho, porém sem garantias de que está correto.
\end{frame}

\begin{frame}{Introdução ao Design de Algoritmos}
    \huge
    \textbf{Exemplo}: Problema de agendamento

    \Large
    \textbf{Entrada}: Conjunto $K$ de $n$ intervalos em uma linha

    \textbf{Saída}: Maior conjunto de intervalos de $K$ que não se sobrepõem
\end{frame}

\begin{frame}[fragile]{Introdução ao Design de Algoritmos}
    \begin{lstlisting}
PrimeiroTrabalho($K$)
  Aceite o primeiro trabalho $t$ de $K$ que não se sobrepõe
  --> a outro trabalho aceito
  Repita até que não sobrem trabalhos
    \end{lstlisting}
\end{frame}

\begin{frame}[fragile]{Introdução ao Design de Algoritmos}
    \begin{lstlisting}
TrabalhoMaisCurto($K$)
  Enquanto $K$ não for vazio faça
    Aceite o trabalho $t$ mais curto em $K$
    Exclua $t$ e qualquer outro trabalho onde há
    --> sobreposição com $t$
    \end{lstlisting}
\end{frame}

\begin{frame}[fragile]{Introdução ao Design de Algoritmos}
    \begin{lstlisting}
AgendamentoExaustivo($K$)
  $j = 0$
  $S_{max} = 0$
  Para cada um dos $2^n$ subconjuntos $S_i$ de intervalos
    Se $S_i$ não tem sobreposições e $tam(S_i) > j$
      então $j = tam(S_i)$ e $S_{max} = S_i$
  Retorne $S_i$
    \end{lstlisting}
\end{frame}

\begin{frame}[fragile]{Introdução ao Design de Algoritmos}
    \begin{lstlisting}
AgendamentoÓtimo($K$)
  Enquanto $K$ não for vazio faça
    Aceite o trabalho $t$ de $K$ que termina mais cedo
    Exclua $t$ e qualquer outro trabalho onde há
    --> sobreposição com $t$
    \end{lstlisting}
\end{frame}

\begin{frame}{Introdução ao Design de Algoritmos}
    \huge
    Algoritmos aparentemente razoáveis podem estar incorretos. A correção de um algoritmo é uma propriedade que precisa ser demonstrada \textbf{cuidadosamente}.
\end{frame}

\begin{frame}{Introdução ao Design de Algoritmos}
    \Large
    Algoritmos podem ser expressados através de:

    \begin{itemize}
        \item \Large Texto
        \item \Large Pseudocódigo
        \item \Large Programa real
    \end{itemize}

    \onslide<2->{\huge O coração do algoritmo é uma \textbf{idéia}.}
\end{frame}

\begin{frame}{Introdução ao Design de Algoritmos}
    \Large
    Especificação de problemas tem duas partes:

    \begin{enumerate}
        \item \Large O conjunto permitido de instâncias de entrada
        \item \Large A propriedades exigidas na saída
    \end{enumerate}
\end{frame}

\begin{frame}{Introdução ao Design de Algoritmos}
    \huge
    Outro lado da moeda: Demonstrando \textbf{incorreção} com contra-exemplos
\end{frame}

\begin{frame}{Introdução ao Design de Algoritmos}
    \Large Propriedades importantes de contra-exemplos:

    \begin{itemize}
        \item \Large Verificabilidade
        \begin{itemize}
            \Large
            \item Calcular a resposta
            \item Mostrar uma resposta melhor
        \end{itemize}
        \item \Large Simplicidade
    \end{itemize}
\end{frame}

\begin{frame}{Introdução ao Design de Algoritmos}
    \Large Algumas estratégias:

    \begin{itemize}
        \Large
        \item Pense pequeno
        \item Pense exaustivamente
        \item Procure extremos (grande/pequeno, esquerda/direita, poucos/muitos, perto/longe)
    \end{itemize}
\end{frame}

\end{document}
